Indeed, in (i), the first assertion follows from 6.7 (i); the second, from 6.10 and 1.8 (i), and from 6.2 (applied to @@M0@@, @@M1@@) and 1.3. Finally, (ii) follows from (i) and 1.13.

\sgaStableTarget{sga6.e13.rem.6-12}\textbf{Remark 6.12.} Let @@M2@@ be an algebraically closed field, and @@M3@@ a projective @@M4@@-scheme of pure dimension @@M5@@, with no embedded associated prime cycle. If we allow arbitrary nilpotent elements in @@M6@@, we can no longer bound the coefficients in the Hilbert polynomial @@M7@@ in terms of @@M8@@ alone.

For example, let @@M9@@ be a smooth curve of degree @@M10@@ in @@M11@@. For every @@M12@@, there exists a smooth surface @@M13@@ of degree @@M14@@ which contains @@M15@@. Then @@M16@@ is a Cartier divisor on @@M17@@; put @@M18@@. Finally, we have the exact sequence @@M19@@ and the formula @@M20@@; hence @@M21@@ (@@M22@@), although @@M23@@.

\sgaStableTarget{sga6.e13.cor.6-13}\textbf{Corollary 6.13.} \emph{Under the conditions of 6.1, consider an exact sequence}
@@M24@@
\emph{with @@M25@@. Then, for @@M26@@,}
@@M27@@

Indeed, put @@M28@@; clearly, @@M29@@ is a @@M30@@-sheaf, and @@M31@@; whence the assertion by 6.7 (i), and the formula @@M32@@.

